% from ega4-2-fr.tex
\oldpage[IV-2]{130}

immergés (5.7.7). L'idéal $\mathfrak p^{(1)}$ est donc le seul idéal premier
associé à $A^{(1)}/tA^{(1)}$, autrement dit, $tA^{(1)}$ est un idéal
$\mathfrak p^{(1)}$-primaire de $A^{(1)}$; cela veut dire encore que
$tA^{(1)}$ est l'image réciproque dans $A^{(1)}$ de
$tA^{(1)}_{\mathfrak p^{(1)}}=tA_{\mathfrak p}$, et ce dernier est l'idéal
maximal de $A_{\mathfrak p}$ d'après (ii); donc
$tA^{(1)}=\mathfrak p^{(1)}$ est premier dans $A^{(1)}$. D'autre part,
$A^{(1)}/\mathfrak p^{(1)}$ est fini sur $A/\mathfrak p$ et a même corps des
fractions (savoir le corps résiduel de
$A^{(1)}_{\mathfrak p^{(1)}}=A_{\mathfrak p}$), donc il est identique à
$A/\mathfrak p$ en vertu de l'hypothèse (iii). On peut donc écrire
$A^{(1)}=A+\mathfrak p^{(1)}=A+tA^{(1)}$, et comme $t$ est contenu dans le
radical de l'anneau semi-local noethérien $A^{(1)}$, le lemme de Nakayama
entraîne $A^{(1)}=A$, d'où $\mathfrak p^{(1)}=\mathfrak p=tA$. Mais comme $A$
est caténaire et $A/tA$ intègre et intégralement clos en vertu de l'hypothèse
(iii), on conclut de (5.12.7) que $A$ est intégralement clos. C.Q.F.D.

\begin{remark}[5.12.9]
\label{IV.2.s5.12.9-fr}
La conclusion de (5.12.8) tombe en défaut si l'on ne suppose plus $A$
équidimensionnel. Par exemple, soient $k$ un corps, $B$ l'anneau de polynômes
$k[X,Y,Z]$ à trois indéterminées, $C$ l'anneau
$B/\mathfrak r_1\mathfrak r_2$, où $\mathfrak r_1=BZ$ et
$\mathfrak r_2=BX+BY$; soient $\mathfrak m$ l'idéal maximal $BX+BY+BZ$,
$\mathfrak n=\mathfrak m/\mathfrak r_1\mathfrak r_2$ son image dans $C$, $A$
l'anneau local $C_{\mathfrak n}$. Enfin, soit $t_0$ l'image dans $C$ de
l'élément $Y+Z$ de $B$, $t$ son image dans $A$
($\operatorname{Spec}(C)$ est donc formé d'un plan $P$ et d'une droite $D$ non
contenue dans ce plan, $\operatorname{Spec}(C/t_0C)$ d'une droite $D'$ contenue
dans $P$ et passant par le point $D\cap P$). L'anneau $A$ est réduit et
caténaire (étant quotient d'un anneau régulier) et ses idéaux premiers minimaux
$\mathfrak q_1$, $\mathfrak q_2$ sont les images de $\mathfrak r_1$,
$\mathfrak r_2$. Le $A$-module $A/tA$ n'a qu'un seul idéal premier associé non
immergé $\mathfrak p$, image de $BY+BZ$, $\mathfrak pA_{\mathfrak p}$ est
engendré par $t/1$ et $A/\mathfrak p$ isomorphe à $k[X]$; mais $A$ n'est pas
intègre.
\end{remark}

\begin{corollary}[5.12.10]
\label{IV.2.s5.12.10-fr}
Soit $A$ un anneau local intègre noethérien, vérifiant l'une des hypothèses
suivantes:
\begin{enumerate}
\item[a)] $A$ est quotient d'un anneau régulier.
\item[b)] $A$ est universellement caténaire et universellement japonais.
\end{enumerate}
Soit $(x_i)_{1\leq i\leq n}$ une suite de $n$ éléments de $A$; posons
$\mathfrak J=x_1A+\cdots+x_nA$, et supposons vérifiées les conditions
suivantes:
\begin{enumerate}
\item[(i)] Il n'existe qu'un seul idéal premier non immergé $\mathfrak p$
associé à $A/\mathfrak J$ et $\mathfrak p$ est de hauteur $n$.
\item[(ii)] L'idéal maximal $\mathfrak pA_{\mathfrak p}$ de $A_{\mathfrak p}$
est engendré par les $x_i/1$ (donc $A_{\mathfrak p}$ est un anneau local
régulier de dimension $n$).
\item[(iii)] L'anneau $A/\mathfrak p$ est intégralement clos.
\end{enumerate}
Sous ces conditions, pour tout entier $i$ tel que $0\leq i\leq n$, l'anneau
quotient $A/(x_1A+\cdots+x_iA)$ est intégralement clos et de dimension
$\dim(A)-i$. En particulier, $\mathfrak J$ est premier, égal à $\mathfrak p$,
$A$ est intégralement clos et $(x_i)_{1\leq i\leq n}$ est une suite
$A$-régulière.
\end{corollary}

Procédons par récurrence sur $n$, la proposition étant triviale pour $n=0$ et
se réduisant au lemme de Hironaka (5.12.8) pour $n=1$. On peut donc supposer
$n\geq2$. Soit $\mathfrak J'=x_1A+\cdots+x_{n-1}A$, et soit $\mathfrak q$ un
idéal premier minimal parmi ceux qui contiennent $\mathfrak J'$; on a
$\operatorname{ht}(\mathfrak q)\leq n-1$ (0, 16.3.1); montrons que l'on a
$\mathfrak q\subset\mathfrak p$. En effet, si $\mathfrak p'$ est minimal parmi
les idéaux premiers de $A$ contenant $\mathfrak q+x_nA$,
$\mathfrak p'/\mathfrak q$ est de hauteur $\leq1$ dans $A/\mathfrak q$ par le
Hauptidealsatz (0, 16.3.2), donc $\mathfrak p'$ est de hauteur $\leq n$ puisque
$A$ est caténaire (0, 16.1.4); mais comme il contient $\mathfrak J$, il est
nécessairement égal à $\mathfrak p$, d'où
$\mathfrak q\subset\mathfrak p$,

